% This section will focus on the uncertainty in localization, kidnapping, and research questions.
% This section will begin with a clear description of uncertainty during localization in robotics.
% It will introduce the concept of kidnapping and provide examples of how this phenomenon can occur in real examples \cite{DBLP:journals/corr/abs-1106-0222}. 
% Next, this section will very briefly introduce, but not explain in detail, the localization algorithms including dead reckoning \cite{ros2_control_mobile_robot_kinematics}, extended kalman filter (EKF) \cite{leonard1991mobile}, and particle filter (PF) \cite{doucet2000sequential}.
% At the end of this section, this section will state the hypotheses that the PF will recover localization accuracy following a kidnapping event more reliably than the EKF and that given a uniform initial pose prior, the EKF will achieve localization accuracy comparable to or better than the PF with lower computational cost.

\section{Problem Statement}\label{sec:Problem-Statement}

In mobile robotics, the planning environment is often projected onto a two dimensional plane.
Let the pose of the robot at time $t$ be defined by the state vector $\mathbf{x}_t = [x_t, y_t, \theta_t]^T$, where $x_t$ and $y_t$ represent the planar coordinates and $\theta_t$ represents the heading. As the robot moves, its true state transitions according to a non-linear motion model:

\begin{equation}
    \mathbf{x}_t = f(\mathbf{x}_{t-1}, \mathbf{u}_t) + \mathbf{w}_t
\end{equation}
where $\mathbf{u}_t$ is the control input (such as translational velocity $v_t$ and rotational velocity $\omega_t$) and $\mathbf{w}_t$ represents process noise caused by factors like wheel slip or mechanical imperfections. 

To correct the drift induced by the process noise, the robot relies on exteroceptive sensors, modeled as:
\begin{equation}
    \mathbf{z}_t = h(\mathbf{x}_t) + \mathbf{v}_t
\end{equation}
where $\mathbf{z}_t$ is the sensor measurement (such as a LiDAR scan) and $\mathbf{v}_t$ represents the measurement noise. 
The core problem of localization is to estimate the \textit{belief}, or the probability distribution over the state $x_t$ at time $t$, conditioned on all past measurements $z_{1:t}$ and all past controls $_{1:t}$\cite{thrun2005probabilistic}.

\begin{equation}\label{eq:belief}
    bel_{x_t} = p(x_t \mid z_{1:t}, u_{1:t})
\end{equation}

\subsection{Kidnapping}
While tracking a robot's pose under normal operating conditions is well understood, catastrophic localization failures remain a significant challenge. 
The most extreme form of this failure is known as the ``kidnapped robot problem'' \cite{DBLP:journals/corr/abs-1106-0222}. 
This occurs when a robot is teleported to a different location without its odometry registering the movement, completely invalidating its current belief state.

In the real world, this could happen if a robot loses power and is re-powered at a different position or if the robot being tracked loses connection with a tracking hub and then reconnects at a different time.
To recover from a kidnapping event, a localization algorithm must be capable of abandoning its incorrect localized belief and search the map to find its new true pose.

\subsection{Bayes Filter}
The most general algorithm for calculating beliefs, shown in Equation \ref{eq:belief}, is the \textbf{Bayes filter}\cite{thrun2005probabilistic}.
The Bayes filter is recursive and calculates the belief $bel(x_t)$ at time $t$ from the $bel(x_{t-1})$ at time $t-1$.
The Bayes filter has two steps: prediction and measurement update.
The prediction calculates the belief $\overline{bel}(x_t)$ over $x_t$ based on the prior belief over $x_{t-1}$ and $u_t$ as shown in Equation \ref{eq:bayes_prediction}.

\begin{equation}\label{eq:bayes_prediction}
    \overline{bel}(x_t)=\int p(x_t \mid u_t, x_{t-1})bel(x_{t-1}) \,dx 
\end{equation}

The measurement update calculates the final belief $bel(x_t)$ by multiplying the belief $\overline{bel}(x_t)$ by the probability that the measurement $z_t$ was observed as shown in Equation \ref{eq:bayes_measurement_update}.

\begin{equation}\label{eq:bayes_measurement_update}
    bel(x_t)= \eta p(z_t \mid x_t) \overline{bel}(x_t) 
\end{equation}

\subsection{Algorithm Overview}
Various algorithms exist to solve the localization problem.
The three algorithms selected for this study have different levels of robustness for kidnapping as well as different expected accuracies and computational costs.
They are all based on the Bayes filter with increasing complexity.

The first algorithm, \textbf{Dead Reckoning}, is the most basic form of state estimation.
It relies only on integrating the control inputs $u_t$ over time to estimate the current pose \cite{ros2_control_mobile_robot_kinematics}.
The Dead Reckoning algorithm implements only the prediction step of the Bayes filter.
Because Dead Reckoning lacks a measurement update step, process noise accumulates with time, making the predicted pose stray from truth with time.
This algorithm is used as a comparative baseline for the other two algorithms.

The second algorithm, \textbf{Extended Kalman Filter (EKF)}, tackles the non-linear motion and measurement models by linearizing them via Taylor series expansion and computing the Jacobians of $f$ and $h$ from a kinematic model and sensor model \cite{thrun2005probabilistic}.
The EKF includes both the prediction step and the measurement update step described by the Bayes filter.
Both Dead Reckoning and EKF are considered \textit{Gaussian Filters} where the probability density over $x_t$ can be characterized by the mean $\mu$ and the covariance $\Sigma$.
Gaussians are unimodal, meaning the posses a single maximum and perform poorly in problems which many distinct hypotheses exist.

The third algorithm, \textbf{Particle Filter (PF)}, is a non-parametric approach that represents the belief state using finite sets of weighted random samples called particles \cite{doucet2000sequential}.
Because the PF does not assume a single Gaussian distribution as the EKF does, the PF can represent complex, multi-modal beliefs.

This study evaluates these algorithms under various factory conditions to test two primary hypotheses. First, it is hypothesized that the Particle Filter will recover localization accuracy following a kidnapping event significantly more reliably than the Extended Kalman Filter. 
This is due to the Particle Filter's ability to maintain multiple pose hypotheses simultaneously. 
Second, it is hypothesized that given a well-defined Gaussian initial pose prior (where the robot's starting location is known), the Extended Kalman Filter will achieve localization accuracy comparable to or better than the Particle Filter while requiring a significantly lower computational cost.
